\documentclass[11pt]{article}

\usepackage{amsmath,amssymb,graphicx}
\usepackage[margin=1in]{geometry}
\usepackage{booktabs}

\title{Emergent Memory Retarded Interaction Dynamics Probe on a Frozen Kernel-Induced Cochain Operator Architecture}
\author{Tomislav Rupi\'c \\ Independent Researcher}
\date{March 2026}

\begin{document}

\maketitle

\begin{abstract}
Stages 10--17 established a falsification-first pre-geometric ladder on a frozen kernel-induced cochain operator architecture. Collective morphology, coarse envelope organization, and localized relational ordering were all observed in restricted forms, but no weak effective binding, relational-feedback, or phase-synchronization channel emerged. Stage 18 changes the coupling class rather than the coupling strength. Instead of testing another instantaneous weak correction, it asks whether a minimal single-lag retarded interaction term can generate collective organization that previous weak channels could not. A narrow Branch-A-first pilot is executed on three canonical representatives: a clustered composite anchor, a counter-propagating corridor, and a phase-ordered symmetric triad. Each case is run under null control, very weak delayed coupling with $\varepsilon=0.01$, and weak delayed coupling with $\varepsilon=0.02$, using one fixed lag tied to the packet width. The result is cleanly negative. All nine runs remain in the label ``no memory dynamic effect''; no composite lifetime gain appears, no basin-persistence gain appears, no new ordering class emerges, and no case is promoted to the $(12 \rightarrow 24)$ refinement ladder. Small changes in delayed-response correlation are detected, but they do not translate into regime-level structural consequences. The bounded conclusion is that, at the tested strengths and lag scale, a single-lag retarded collective term is provisionally morphologically inert on the frozen architecture.
\end{abstract}

\section{Introduction}

The program through Stage 17 is cumulative and deliberately restrictive. Stage 10 established a reproducible morphology instrument for packet transport on the frozen architecture. Stage 11 extended that instrument to collective interaction structure. Stage 12 coarse-grained multi-packet fields into envelope and basin summaries. Stage 13 coupled interaction and coarse structure directly to refinement and showed that coarse organization sharpens under refinement while composite interaction labels remain fragile. Stage 14 then tested several weak structural branches and found no effective binding channel at the tested strengths. Stage 15 asked whether frozen collective data already support a stable relational ordering that behaves metric-like under refinement. The signal was mixed but informative: localized ordering appeared, but not on the strongest composite anchor, and transport-delay ordering remained weak. Stage 16 then asked whether that relational ordering could act back on the dynamics, and the answer was negative. Stage 17 asked the analogous question for weak phase synchronization and again found no emergent collective timing channel.

Stage 18 therefore changes the coupling class rather than the coefficient scale. If the architecture resists new collective law under instantaneous weak feedback, perhaps history dependence is the missing ingredient. The first Stage 18 probe is intentionally narrow. It asks only whether a single-lag retarded self-feedback term can create measurable collective organization on the same frozen architecture.

\section{Frozen Context and Branch-A Pilot Design}

All kernels, projectors, operator constructions, packet families, and coarse-graining procedures remain frozen. The projected transverse sector and periodic boundary condition are retained, together with the packet family inherited from the strongest coherent transverse seed: bandwidth $w=0.1$, cosine carrier, central wave number $k=1.0$, and amplitude scale $0.5$.

The Stage 18A Branch-A-first pilot uses exactly three representatives:
\begin{itemize}
\item clustered composite anchor,
\item counter-propagating corridor,
\item phase-ordered symmetric triad.
\end{itemize}

These span the strongest composite anchor, a transport-style control, and a structured ordering candidate. Each case is run under three conditions:
\begin{itemize}
\item null control,
\item very weak delayed coupling with $\varepsilon=0.01$,
\item weak delayed coupling with $\varepsilon=0.02$.
\end{itemize}

This gives a total of nine base runs. The broader Stage 18 design allows a full 27-run memory-class sweep and a $(12 \rightarrow 24)$ refinement ladder, but the Branch-A-first screen is designed to stop early if the single-lag class already fails. No refinement follow-up is permitted unless one of the nine base runs shows a real morphology or ordering gain.

\section{Single-Lag Retarded Coupling}

The Stage 18A mechanism is deliberately minimal. The instantaneous weak correction
\[
F(t) \mapsto F(t) + \varepsilon G[\mathrm{state}(t)]
\]
is replaced by a retarded term of the form
\[
F(t) \mapsto F(t) + \varepsilon G[\mathrm{state}(t-\tau)].
\]

In the executable pilot, the delayed observable is a local envelope profile sampled from the packet component itself. Each packet carries a short history buffer, and the feedback uses one fixed delay scale tied to the packet width,
\[
\tau / w = 0.5.
\]
This choice is frozen for the Branch-A screen. Delay scanning is deferred unless the first screen shows a real positive signal.

The term is weak, local in packet-component space, and bounded. It does not modify the frozen kernel, projector, or base operator. The purpose is only to test whether a minimal memory-bearing channel can produce collective organization that instantaneous weak channels could not.

\section{Measurement Set and Early-Stop Rule}

Stage 18A reuses the morphology and coarse observables frozen in earlier stages:
\begin{itemize}
\item composite lifetime,
\item binding persistence,
\item overlap evolution,
\item coarse basin persistence.
\end{itemize}

It adds one memory-specific scalar diagnostic:
\begin{itemize}
\item delayed-response correlation.
\end{itemize}

For each delayed run, four deltas are computed relative to the null control of the same representative:
\begin{itemize}
\item composite lifetime delta,
\item binding persistence delta,
\item coarse persistence delta,
\item delayed-response advantage.
\end{itemize}

The allowed output labels are:
\begin{itemize}
\item no memory dynamic effect,
\item transient history response,
\item scale-fragile memory-assisted ordering,
\item scale-stable memory-assisted ordering.
\end{itemize}

The early-stop rule is strict. Stage 18A freezes negative if all nine base runs show:
\begin{itemize}
\item no composite gain,
\item no basin-persistence gain,
\item no new ordering class,
\item no meaningful delayed-response advantage that crosses the promotion gate together with morphology change.
\end{itemize}

That rule is exactly what occurs.

\section{Pilot Outcome}

The Stage 18A result is uniform:
\begin{itemize}
\item no memory dynamic effect: 9,
\item transient history response: 0,
\item scale-fragile memory-assisted ordering: 0,
\item scale-stable memory-assisted ordering: 0.
\end{itemize}

No case is promoted to the $(12 \rightarrow 24)$ refinement ladder.

\subsection{Clustered composite anchor}

The clustered composite anchor remains in the \emph{metastable composite regime} with \emph{diffuse coarse field regime} under all three conditions. Composite lifetime remains fixed at $1.4764$, binding persistence remains fixed at $1.0000$, and coarse persistence remains fixed at $1.0000$. The delayed-response advantage increases slightly from null to $0.0123$ at $\varepsilon=0.01$ and to $0.0243$ at $\varepsilon=0.02$, but those shifts do not produce any regime change or promotion.

\subsection{Counter-propagating corridor}

The corridor remains in the \emph{transient binding regime} with \emph{diffuse coarse field regime} under all three conditions. Composite lifetime remains $0.9842$, binding persistence remains $0.6667$, and coarse persistence remains $1.0000$. The delayed-response advantage changes by only $2.1\times 10^{-4}$ at $\varepsilon=0.01$ and $4.2\times 10^{-4}$ at $\varepsilon=0.02$.

\subsection{Phase-ordered symmetric triad}

The symmetric triad remains in the \emph{transient binding regime} with \emph{diffuse coarse field regime} under all three conditions. Composite lifetime and binding persistence remain at zero, while coarse persistence remains $1.0000$. The delayed-response advantage changes modestly, reaching $0.0102$ at $\varepsilon=0.01$ and $0.0206$ at $\varepsilon=0.02$, but again without any morphology consequence.

\section{Interpretation of the Negative Result}

The Stage 18A null is useful for three reasons.

First, it shows that the frozen architecture is not secretly waiting for a trivial delay effect to unlock a new collective regime. A physically motivated change in coupling class was introduced, and the existing basins remained intact.

Second, it shows that memory is not automatically an ordering resource in this architecture. The delayed-response correlation shifts indicate that the delayed term is not entirely invisible. The system can register the lag statistically. But it does not convert that signal into morphology gain, basin reshaping, or a new ordering class.

Third, it validates the Branch-A-first early-stop logic. The pilot was designed precisely to prevent unnecessary expansion into a full 27-run memory-class sweep unless a real gain appeared. No such gain appeared, so the negative screen itself is the correct scientific product.

The clean bounded statement is therefore:

\begin{quote}
A single-lag retarded collective coupling term produces measurable delayed-response correlation shifts but no detectable change in dynamical regime structure, basin persistence, or composite stability at the tested strengths.
\end{quote}

\section{Interpretation Boundary}

Stage 18A does not support:
\begin{itemize}
\item effective binding claims,
\item spacetime claims,
\item causal-law claims,
\item memory-induced geometry claims.
\end{itemize}

The only bounded claims supported here are:
\begin{itemize}
\item a single-lag retarded term can be implemented cleanly on the frozen architecture,
\item at $\varepsilon=0.01$ and $0.02$ with $\tau / w = 0.5$ it produces small delayed-response correlation shifts,
\item those shifts have no measurable morphology, basin, or ordering consequence in the tested representatives,
\item the mechanism is provisionally morphologically inert at the tested scale.
\end{itemize}

\section{Conclusion}

Stage 18A is the first direct test of whether a minimal retarded collective term can create a new organizing channel on the frozen architecture. The answer, in the nine-run Branch-A-first pilot, is negative. Across a clustered composite anchor, a counter-propagating corridor, and a phase-ordered symmetric triad, weak delayed coupling produces no detectable regime-level consequence.

This result tightens the map rather than closing it. By this point the program has tested weak binding, weak relational feedback, weak phase synchronization, and now weak single-lag memory. None has produced a new collective law. The current architecture remains capable of structured morphology and localized ordering, but it still resists promotion of that structure into an emergent organizing dynamics at weak coupling.

The next sensible move, if memory space is revisited at all, is not stronger $\tau$ or stronger $\varepsilon$ on the same branch. It is a qualitatively different memory topology: distributed memory kernels, state-dependent lag, or multiscale history weighting.

\end{document}
